% 图契约：本题的源、边界和推导逐项对应；箭头符号按图注解释。
\documentclass[tikz,border=7pt]{standalone}
\usepackage{fontspec,xeCJK,amsmath,amssymb}
\IfFontExistsTF{Times New Roman}{\setmainfont{Times New Roman}}{\setmainfont{TeX Gyre Termes}}
\IfFontExistsTF{Heiti SC}{\setCJKmainfont[AutoFakeBold=1.5]{Heiti SC}}{\setCJKmainfont[AutoFakeBold=1.5]{FandolHei-Regular}}
\usetikzlibrary{arrows.meta,bending,calc,patterns,positioning}
\definecolor{ink}{HTML}{243745}
\definecolor{blue}{HTML}{275DAB}
\definecolor{green}{HTML}{177C63}
\definecolor{red}{HTML}{BD4935}
\definecolor{orange}{HTML}{AA6C16}
\definecolor{purple}{HTML}{715393}
\tikzset{every node/.style={font=\fontsize{9.3}{12.5}\selectfont,text=ink},
axis/.style={-{Stealth[length=4pt]},line width=.65pt,draw=ink!65},
vec/.style={-{Stealth[length=5pt]},line width=1pt,draw=blue},
short/.style={-{Stealth[length=3pt]},line width=.8pt},
guide/.style={dashed,draw=ink!40,line width=.5pt},
card/.style={rounded corners=2pt,draw=blue!65,fill=blue!4,align=center,text width=6.5cm,inner xsep=9pt,inner ysep=7pt},
title/.style={font=\fontsize{10.5}{14}\selectfont\bfseries,text width=15cm,align=center},
label/.style={font=\fontsize{8.4}{11}\selectfont,align=center},
note/.style={label,text width=6.2cm}}
\begin{document}
\begin{tikzpicture}[x=1cm,y=1cm]
\path[use as bounding box] (-8,-3.5) rectangle(8,3.5);
\node[title] at (0,3.15) {M2-017\quad 远场展开中两项单极贡献抵消};
\begin{scope}[shift={(-4,0)}]
\draw[axis] (0,-1.9)--(0,2.1) node[above] {$z$};
\fill[red] (0,0.9) circle(2.3pt) node[right=3pt] {$+q$};
\fill[blue] (0,-0.9) circle(2.3pt) node[right=3pt] {$-q$};
\draw[vec,green] (.9,-.65)--(.9,.65) node[midway,right] {$\mathbf p$};\node[left] at (0,.9) {$a$};\node[left] at (0,-.9) {$-a$};
\node[label] at(1.65,1.65){$P:r\gg a$};\node[note] at (0,-2.85) {示意观测点在偶极正轴；有限电荷对仍有高阶修正。};
\end{scope}
\node[card,draw=blue!65,fill=blue!4] at (3.9,1.8) {$\displaystyle\frac1{r\pm a}\simeq\frac1r\mp\frac a{r^2}$};
\node[card,draw=green!65,fill=green!4] at (3.9,0) {$\displaystyle\frac1{r-a}-\frac1{r+a}\simeq\frac{2a}{r^2}$};
\node[card,draw=orange!65,fill=orange!4] at (3.9,-1.8) {两个 $r^{-1}$ 项相消；\\剩下偶极阶 $r^{-2}$。};
\end{tikzpicture}\end{document}